\section*{Capítulo \ref{ch:g1:light-and-shadows} --- Luz e sombras}

\begin{solution}{exo:g1:light-and-shadows:1}
Fontes: a chama da vela, o vaga-lume, a lanterna. Não fontes: a Lua,
o espelho, a parede branca --- eles só devolvem a \omterm{def:g1:light-and-shadows:source}{luz} que cai sobre
eles.
\end{solution}

\begin{solution}{exo:g1:light-and-shadows:2}
Falta \omterm{def:g1:light-and-shadows:source}{luz}. Não há fonte nenhuma no cômodo, então nenhuma \omterm{def:g1:light-and-shadows:source}{luz} cai
sobre as coisas, e coisas sem \omterm{def:g1:light-and-shadows:source}{luz} sobre elas não podem ser vistas.
\end{solution}

\begin{solution}{exo:g1:light-and-shadows:3}
Uma \omterm{def:g1:light-and-shadows:source}{fonte de luz}, um obstáculo e uma parede (ou tela, ou o chão) ---
nessa ordem, ao longo de uma linha reta: a fonte acende, o obstáculo
barra, a parede mostra a mancha escura.
\end{solution}

\begin{solution}{exo:g1:light-and-shadows:4}
A \omterm{def:g1:light-and-shadows:shadow}{sombra} fica à sua direita --- do lado oposto ao Sol. Depois da
meia-volta, o Sol fica à sua direita, então a \omterm{def:g1:light-and-shadows:shadow}{sombra} se deita à sua
esquerda: sempre do lado contrário ao do Sol.
\end{solution}

\begin{solution}{exo:g1:light-and-shadows:5}
Levar o brinquedo \emph{para perto da lâmpada} deixa a \omterm{def:g1:light-and-shadows:shadow}{sombra}
enorme; levá-lo \emph{para perto da parede} encolhe a \omterm{def:g1:light-and-shadows:shadow}{sombra} até o
tamanho do brinquedo.
\end{solution}

\begin{solution}{exo:g1:light-and-shadows:6}
Não. Uma \omterm{def:g1:light-and-shadows:shadow}{sombra} não é uma coisa feita de alguma matéria --- ela é um
lugar onde falta \omterm{def:g1:light-and-shadows:source}{luz} porque um obstáculo mantém a \omterm{def:g1:light-and-shadows:source}{luz} fora dali. Como
não há nada ali, não há nada para pegar.
\end{solution}

\begin{solution}{exo:g1:light-and-shadows:7}
O Sol é uma fonte muito mais forte do que qualquer lâmpada: a \omterm{def:g1:light-and-shadows:source}{luz}
dele é tão forte que pode machucar os olhos para sempre, e não apenas
ofuscá-los por um instante.
\end{solution}

\begin{solution}{exo:g1:light-and-shadows:8}
No fim da tarde, quando o Sol fica baixo. O desenho deve mostrar uma
\omterm{def:g1:light-and-shadows:shadow}{sombra} curta ao pé da árvore ao meio-dia e uma \omterm{def:g1:light-and-shadows:shadow}{sombra} comprida e
esticada à tarde, as duas apontando para longe do Sol.
\end{solution}

\begin{solution}{exo:g1:light-and-shadows:9}
Sam está certo. A \omterm{def:g1:light-and-shadows:shadow}{sombra} ficou o tempo todo exatamente do lado oposto
ao Sol. Mia é que girou, então a \omterm{def:g1:light-and-shadows:shadow}{sombra} mudou de lugar \emph{em
relação à frente e às costas dela} --- mas no chão a \omterm{def:g1:light-and-shadows:shadow}{sombra} não se
mexeu.
\end{solution}
